% 第 12 题：电阻丝（圆形表盘刻度）+ 实物接线示意图（按原图矢量重绘）
\begin{tikzpicture}[font=\small, >=Stealth, scale=0.9]
  % 圆形电阻丝（表盘）
  \draw[thick] (5.4,4.05) circle (2.15);
  \node at (5.4,4.05) {表盘};
  % 刻度与数字（11 到 5 环绕）
  \foreach \k in {0,1,...,11} {
    \pgfmathsetmacro\ang{102-30*\k}
    \pgfmathsetmacro\lab{\k+1}
    \draw (5.4,4.05) ++(\ang:2.15) -- ++(\ang:-0.22);
    \node[font=\scriptsize] at ([shift={(\ang:1.78)}]5.4,4.05) {\lab};
  }
  % 端点 a、b
  \fill (4.05,2.35) circle (2.6pt);
  \node[below left] at (4.05,2.35) {$a$};
  \fill (6.75,2.35) circle (2.6pt);
  \node[below right] at (6.75,2.35) {$b$};
  % 电阻丝标注
  \draw[->] (2.0,4.55) -- (4.3,5.55);
  \node[left] at (1.95,4.55) {电阻丝};
  % 微安表
  \draw[thick] (0.45,2.05) .. controls (0.35,3.15) and (2.05,3.15) .. (1.95,2.05)
      -- (0.45,2.05);
  \draw[thick] (0.65,2.55) rectangle (1.75,2.95);
  \draw[thick] (0.75,2.62) -- (1.15,2.86);
  \node[font=\tiny] at (1.35,2.78) {$\mu$A};
  \node (Am) at (1.2,2.3) {};
  % R0（电阻箱/滑动变阻器）
  \draw[thick] (8.9,3.55) rectangle (11.4,4.35);
  \draw[thick] (9.6,4.35) -- (9.6,4.95) -- (10.7,4.95) -- (10.7,4.35);
  \fill[gray!35] (9.75,4.6) rectangle (10.55,4.95);
  \node[above] at (10.15,4.95) {$R_{0}$};
  % 金属夹
  \draw[thick] (8.55,4.35) -- (8.7,4.9) node[above,font=\scriptsize] {};
  \node[above,font=\scriptsize] at (8.3,5.15) {金属夹};
  % 电池组
  \draw[thick] (2.15,0.35) rectangle (4.35,1.35);
  \foreach \x in {2.35,2.75,3.15,3.55} { \draw[thick] (\x,0.45) -- (\x,1.25); }
  \node[above,font=\scriptsize] at (2.5,1.35) {$+$};
  \node[above,font=\scriptsize] at (4.0,1.35) {$-$};
  % 开关
  \draw (8.1,0.35) circle (0.06);
  \draw (9.4,0.35) circle (0.06);
  \draw[thick] (8.1,0.42) -- (9.25,0.85);
  \draw (8.1,0.35) -- (8.1,0.1) -- (9.4,0.1) -- (9.4,0.35);
  % 连线
  \draw[thick] (0.55,2.05) -- (0.3,2.05) -- (0.3,0.85) -- (2.15,0.85);
  \draw[thick] (1.85,2.05) -- (4.05,2.35);
  \draw[thick] (6.75,2.35) -- (11.9,2.35) -- (11.9,3.95) -- (11.4,3.95);
  \draw[thick] (8.9,3.95) -- (7.3,3.95) -- (7.3,1.25);
  \draw[thick] (7.3,1.25) -- (7.3,0.35) -- (8.1,0.35);
  \draw[thick] (9.4,0.1) -- (9.4,-0.15) -- (4.35,-0.15) -- (4.35,0.35);
\end{tikzpicture}
